\errorcontextlines9
\newif\ifshowmacros
%\showmacrostrue

\ifx\usepackage\RequirePackage\else
\documentclass{article}
\usepackage{float}
\fi

\makeatletter
%%%      \gdef\@chapter[#1]#2{\relax\relax\relax\@empty}%
\makeatother

\usepackage[rightFloats]{fltpage}
\usepackage{rotating}
\usepackage{sidecap}
\usepackage{wrapfig}
\usepackage{listings}

%%%\usepackage{hyperref}
\usepackage[figurename=Abbildung,listfigurename={Liste der Abbildungen},within=none]{newfloat}
\ifdefined\chapter
  \newfloatsetup{chapterlistsgap=20pt}
\fi
\SetupFloatingEnvironment{figure}{within=section}
\ForEachFloatingEnvironment{\typeout{#1}}
\ForEachFloatingEnvironment*{\typeout{*#1}}

\makeatletter

\DeclareFloatingEnvironment{diagram}
\ifcsname theHdiagram\endcsname
% ok
%\show\theHdiagram
\else
% wrong
\PackageError{newfloat}{theHdiagram should be defined!}{}
\endinput
\fi
\ifdefined\chapter
  \SetupFloatingEnvironment{diagram}{within=chapter}
\fi
\ifshowmacros
\showthe\c@float@type
\show\thediagram
\show\theHdiagram
\show\fps@diagram
\show\ftype@diagram
\show\ext@diagram
\show\fnum@diagram
\show\listofdiagrams
\show\l@diagram
\show\listdiagramname
\show\diagramname
\fi

\DeclareFloatingEnvironment[placement=h,fileext=lop,within=section]{program}
  [Programm][Programmverzeichnis]
\ifdefined\chapter
  \SetupFloatingEnvironment{program}{chapterlistsgaps=off}
\fi
\ifshowmacros
\showthe\c@float@type
\show\theprogram
\show\theHprogram
\show\fps@program
\show\ftype@program
\show\ext@program
\show\fnum@program
\show\listofprograms
\show\l@program
\show\listprogramname
\show\programname
\fi

\ifdefined\chapter
  \SetupFloatingEnvironment{lstlisting}{chapterlistsgaps=on}
\fi

\ifshowmacros
\show\newfloat@typelist
\fi

\makeatother

\begin{document}
\ifshowmacros
\show\FPdiagram
\show\endFPdiagram
\show\SCdiagram
\show\endSCdiagram
\fi

\tableofcontents
\listoffigures
\listofdiagrams
\listofprograms
\lstlistoflistings

\ifdefined\chapter \else
\def\chapter{\clearpage\origsection}
\let\origsection\section
\let\section\subsection
\let\subsection\subsubsection
\fi

%\chapter{plain}
\clearpage

\begin{figure}[!ht]
\centering A
\caption{A figure}
\end{figure}

\begin{diagram}[!ht]
\centering xyz
\caption{A diagram}
\label{diagram}
\end{diagram}

\begin{program}[!ht]
  \verb|#include <stdio.h>|\\
  \verb|int main()|\\
  \verb|{|\\
  \verb|    printf( "Hello world!\n" );|\\
  \verb|    return 0;|\\
  \verb|}|
  \caption{A very simple C program}
\end{program}

   \begin{lstlisting}[caption={Useless code},label=useless1,captionpos=t]
   for i:=maxint to 0 do
   begin
       { do nothing }
   end;
   \end{lstlisting}

\ifdefined\SCdiagram
\chapter{fltpage}

The following example will produce a caption below and an almost empty facing
page:
\begin{FPdiagram}
 \caption[A caption alone on a page]%
         {A caption alone on a page will belong to the
          following figure without a caption!}
 \newlength\tempdima
 \setlength\tempdima\textheight
 \addtolength\tempdima{-7pt}%
 \fbox{\rule[-10cm]{0pt}{\tempdima}Frame symbolizing figure on a whole page.\hspace{5cm}}
\end{FPdiagram}
\begin{quote}
  \verb|\usepackage[rightFloats]{fltpage}|\\
  \ldots\\
  \verb|\begin{FPdiagram}|\\
  \verb|  \caption{A caption alone ... figure without a caption!}|\\
  \verb|  \fbox{\rule[-10cm]{0pt}{\textheight}Frame ... page.\hspace{5cm}}|\\
  \verb|\end{FPdiagram}|
\end{quote}
\fi

\ifdefined\sidewaysdiagram
\chapter{rotating}

\begin{sidewaysdiagram}
\centering
\strut Oh wie sch\"on ist doch der Westerwald!
\caption{A sidewaysdiagram}
\end{sidewaysdiagram}
\begin{sidewaysprogram}
  \verb|#include <stdio.h>|\\
  \verb|int main()|\\
  \verb|{|\\
  \verb|    printf( "Hello world!\n" );|\\
  \verb|    return 0;|\\
  \verb|}|
  \caption{A very simple C program}
\end{sidewaysprogram}
\fi

\ifdefined\SCdiagram
\chapter{sidecap}

\begin{SCdiagram}
\strut Oh wie sch\"on ist doch der Westerwald!
\caption{A SCdiagram}
\end{SCdiagram}
\begin{SCprogram}
  \verb|#include <stdio.h>|\\
  \verb|int main()|\\
  \verb|{|\\
  \verb|    printf( "Hello world!\n" );|\\
  \verb|    return 0;|\\
  \verb|}|
  \caption{A very simple C program}
\end{SCprogram}
\fi

\ifdefined\wrapdiagram
\chapter{wrapfig}

\newcommand\sample{Some text for our page
 that is reused over and over again. }

The starting place for the wrapfigure
environment was manually determined in
the current exsample by first setting the text without
\begin{wrapfigure}[7]{l}[0.2\width]{0pt}
  \centering
  \fbox{This is a ``wrapfigure''.}
  \caption{An example of the
    \texttt{wrapfigure} environment}
\end{wrapfigure}
the figure to find the linebreaks.

\sample \sample \sample
\sample \sample \sample
\sample \sample \sample

The starting place for the wrapfigure
environment was manually determined in
the current exsample by first setting the text without
\begin{wrapdiagram}[7]{l}[0.2\width]{0pt}
  \centering
  \fbox{This is a ``wrapdiagram''.}
  \caption{An example of the
    \texttt{wrapdiagram} environment}
\end{wrapdiagram}
the figure to find the linebreaks.

\sample \sample \sample
\sample \sample \sample
\sample \sample \sample
\fi

\chapter{Appendix}

   \begin{lstlisting}[caption={Useless code},label=useless2,captionpos=t]
   for i:=maxint to 0 do
   begin
       { do nothing }
   end;
   \end{lstlisting}

\end{document}
